% ============================================================
% Chapter 5: Development
%   Section: Technology Stack
% Included from chapters/05-development/development.tex
% Sources: docs/arquitectura-tecnica.md 2.1; docs/decisions/inference.md;
% docs/roadmap.md (the llama3.1 -> llama3.2:3b swap, Sep 2026).
% ============================================================

\section{Technology Stack}
\label{sec:stack}

\begin{table}[ht]
\centering
\small
\caption{The technology stack and why each piece was chosen.}
\label{tab:stack}
\begin{tabularx}{\linewidth}{@{}lXX@{}}
\toprule
\textbf{Layer} & \textbf{Technology} & \textbf{Why} \\
\midrule
Backend & Python 3.12, FastAPI, \texttt{uv} & The NLP and scraping
libraries are Python; FastAPI gives typed request models and
background jobs; \texttt{uv} locks dependencies reproducibly \\
Frontend & Next.js 14 (App Router), TypeScript & Server-side API
proxies keep the API key out of the browser; a hand-written CSS design
system, no UI library \\
Services & Docker Compose & The same stack in development, evaluation
and production \\
Search & SearXNG (self-hosted) & Free, no key, and a closed list of
engines chosen and measured one by one \\
Extraction & trafilatura \cite{barbaresi2021trafilatura},
BeautifulSoup, Playwright (optional) & A cascade from cheapest to most
expensive (Subsection~\ref{subsec:scraping-cascade}) \\
Entities & GLiNER (\texttt{gliner\_small})
\cite{zaratiana2024gliner} & Entity types given at run time, no
training; the smallest of the three models \\
Sentiment & \texttt{twitter-xlm-roberta-}\allowbreak\texttt{base-sentiment}
\cite{barbieri2022xlmt} & Multilingual, covers English and Spanish
with one model \\
Embeddings & BGE-M3 \cite{chen2024bgem3} & Multilingual, 1,024
dimensions; one space for every similarity in the system \\
Keywords & YAKE \cite{campos2020yake} & Unsupervised, per-language
stopwords \\
Language model & \texttt{llama3.2:3b} through Ollama & About 2~GB,
fits the development GPU; OpenAI-compatible, so replaceable by
configuration \\
Vector store & Qdrant (embedded) & No server to run; one collection of
article embeddings \\
Graph & Neo4j Community & A property graph with a query language;
read-only console for exploration \\
Infrastructure & Terraform, Google Cloud, Caddy & One virtual machine
running the same compose stack; HTTPS and a site password in front \\
\bottomrule
\end{tabularx}
\end{table}

\paragraph{The language model.} The default model was
\texttt{llama3.1} (8 billion parameters, about 4.9~GB) until September,
when it was replaced by \texttt{llama3.2:3b} (about 2~GB). The change
was made for resources --- the smaller model fits the development GPU
and a CPU-only virtual machine --- not because it was measured to be
better. Measuring verifier models is the purpose of the model
benchmarks (Section~\ref{sec:benchmarks}); until then the choice is a
resource decision and is reported as such.

\paragraph{No linter.} The repository has no linter or formatter
configured; the decision was recorded rather than left implicit. The
rules that are costly to break are enforced as tests instead
(Section~\ref{sec:dev-testing}).
